\section{Background}

Our algorithm is motivated by a classic question in numerical linear algebra.

\subsection{Problem}

Our goal is, given a matrix $ X \in \mathbb R^{D \times P}$, to select a subset $ S \subset [P]$ with $| S| = D$ such that $X_{.  S}$ is as orthonormal as possible in a computationally efficient way.
To this end, we define a ground truth loss function that measures orthonormalness, and then introduce a surrogate loss function that convexifies the problem so that it may be efficiently solved.
Figure \ref{fig:data} and \ref{fig:selection} provide a visual depiction of this problem when $D=2$ and $P=30$.

\subsection{Isometry}

Isometries are a classical object in numerical linear algebra and geometry.
The mathematical construct central to our estimation problem is the isometric linear map.
\begin{definition}
\label{def:isometry}
A matrix $X \in \mathbb R^{D \times P}$ is an \textbf{isometry} or, equivalently, \textbf{unitary} if
\begin{align}
X X^T = I_D.
\end{align}
\end{definition}
Note that necessarily, $P \geq D$.

The analogous object for smooth nonlinear maps is defined through the differential.
\begin{definition}
The \textbf{differential} of a smooth map $\phi:\mathcal M \to \mathcal N$ between $D$ dimensional manifolds $\M \subseteq \mathbb R^B$ and $\N \subseteq \mathbb R^P$ is a map in tangent bases $x_1 \dots x_{D}$ of $T_\xi \M$ and $y_1 \dots y_{D}$ of $T_{\phi(\xi)} \N$ consisting of entries
\begin{align}
\label{eq:diff}
    D\phi (\xi) = \begin{bmatrix}
    \frac{\partial \phi_1  }{\partial x_1}(\xi)  & \dots & \frac{\partial \phi_1 }{\partial x_D}(\xi)  \\
    \vdots & & \vdots \\
    \frac{\partial \phi_D }{\partial x_1}(\xi)  & \dots & \frac{\partial \phi_{D}  }{\partial x_{D}}(\xi) 
    \end{bmatrix}.
\end{align}
\end{definition}

\begin{definition}
\label{def:isometric_at_a_point}
A map $\phi$ between $D$ dimensional submanifolds with inherited Euclidean metric $\mathcal M \subseteq \mathbb R^{B}$ and $\mathcal N  \subseteq \mathbb R^{P}$ is an \textbf{isometry at a point} $\xi \in \mathcal M$ if
\begin{align}
{D \phi (\xi)}^T D \phi (\xi) = I_D.
\end{align}
That is, $\phi$ is an isometry at $\xi$ if $D \phi (\xi)$ is orthonormal.
\end{definition}

\begin{definition}
\label{def:conformal_at_a_point}
A map $\phi$ between $D$ dimensional submanifolds with inherited Euclidean metric $\mathcal M \subseteq \mathbb R^{B}$ and $\mathcal N  \subseteq \mathbb R^{P}$ is \textbf{conformal at a point} $\xi \in \mathcal M$ if there exists a scalar $\lambda(\xi) > 0$ such that
\begin{align}
{D \phi (\xi)}^T D \phi (\xi) = \lambda(\xi) I_D.
\end{align}
That is, $\phi$ is conformal at $\xi$ if its differential preserves angles but may rescale lengths.
\end{definition}

Definition \ref{def:isometry} has several related formulations.
\begin{proposition}
\label{prop:orthonormal_spectrum}
The singular values $\sigma_1(X) \dots \sigma_D(X)$ are equal to $1$ if and only if $X$ is an isometry.
\end{proposition}
\begin{proposition}
\label{prop:orthonormal_basis}
If $P = D$, the component vectors $X_{.1} \dots X_{.D} \in \mathbb R^D$ form an isometry if and only if, for all $d_1, d_2 \in [D], \langle X_{.d_1}, X_{.d_2} \rangle = \begin{cases}
1 \; d_1 = d_2 \\ 
0 \; d_1 \neq d_2 
\end{cases}$.
\end{proposition}

Distance-preserving maps between Riemannian manifolds are smooth isometries \citep[Problem~6-7]{Lee2018-rm}.
This motivates embedding algorithms such as Local Tangent Space Alignment \citep{ZhangZ:04}, Multidimensional Scaling \citep{ChenBuja:localMDS09}, and Isomap \citep{tenenbaum2000ggf} to non-parametrically estimate embeddings that are as close to isometric as possible.

% \begin{proposition}
% \label{prop:distance_differential_isometry}
% Let $\mathcal M$ and $\mathcal N$ be connected Riemannian manifolds with Riemannian distances $d_\mathcal M$ and $d_\mathcal N$.
% A map $\phi:\mathcal M \to \mathcal N$ preserves distances,
% \begin{align}
% d_\mathcal N(\phi(\xi_1), \phi(\xi_2)) = d_\mathcal M(\xi_1, \xi_2),
% \end{align}
% if and only if $\phi$ is smooth and is an isometry at every point \citep[Problem~6-7]{Lee2018-rm}.
% \end{proposition}

\subsection{Coordinate selection via linearization}

More recently, methods for selecting $D$ coordinates from a dictionary $f_{1:P}$ of candidate features have emerged \citep{Kohli2021-lr, Koelle2022-ju, Koelle2024-}.
One motivation for this approach to from a dictionary of candidate features that have been pre-identified as interpretable \citep{Koelle2024-no}. 
In contrast, \citet{Kohli2021-lr} seeks selection of a subset of coordinates that are as conformal as possible, in a manner that we will explain in the next section.
In \citet{Kohli2021-lr}, this loss function is computer for each individual data point, while in \citet{Koelle2022-ju, Koelle2024-}, the loss function is computed over the entire dataset.
However, both of these methods in essence linearize the following best subset selection problem by considering the differentials of the candidate features.

The dictionary-based coordinate search problem best subset selection problem is
\begin{align}
\label{prog:function_best_subset_selection}
\arg \min_{ S \in \binom{[P]}{D}} \ell ( f_{. S})
\end{align}
where $\ell$ is some loss function that measures isometricness of a selected functional subset $f_{. S}$.
In contrast, the linearized version of this problem is
\begin{align}
\label{prog:linearized_best_subset_selection}
\arg \min_{ S \in \binom{[P]}{D}} \ell ( D f_{. S}(\xi))
\end{align}
where $\ell$ is now a loss function that measures the orthonormalness of a matrix.



Due to the relation between metric and global properties of smooth maps, it is possible to linearize the best subset selection problem by considering differentials.

For example, the distortion loss function used the evaluate isometricness in \citep{Kohli2021-lr} reduces in the linear tangent space limit to the condition number.
The \textit{condition number} \citep{GolubWilkinson-condition,Boyd2004-ql} $\frac{\sigma_1(X)}{\sigma_D(X)}$ measures the ratio of the largest and smallest singular values.
This is also known as the \textit{distortion}.
Plainly, this is minimized by conformal maps.

While it is minimized by an isometric function, it is also minimized by any conformal (i.e. low distortion) linear map.

Another example, \citep{Koelle2022-ju}, uses gradients of coordinate functions computed by automatic differentiation.
This is justified by theory about rank conditions, but plainly more specifics are possible.
Also, it considers support recovery over the whole manifold.

\subsection{Linear isometry selection}

A class of functions more suitable for measuring isometry are S-divergences \citep{Sra2011-sdivergence}.
S-divergences define a distance $D(P,P')$ on positive definite matrices $P,P'$ given by 
\begin{align}
D(P,P') = (\log \det (\frac{P + P'}{2}) - \frac{1}{2}\log \det(PP'))^{\frac{1}{2}}.
\end{align}
Applied to the positive definite matrices $X^T X$ and $I_D$, the S-divergence is a log singular value penalty.

\begin{proposition}
\label{prop:sdiv_singular_values}
If $X^T X$ is positive definite, then
\begin{align}
D(X^T X,I_D)^2
&= \sum_{d=1}^D \log \frac{\sigma_d(X) + \sigma_d(X)^{-1}}{2}.
\end{align}
\end{proposition}
By Proposition \ref{prop:orthonormal_spectrum}, this function is uniquely minimized by isometric matrices.
%TODO: explain how works for PSD not just PD.

\subsection{Best subset selection}

Equipped with function $\ell(X)$ such as $D(X,I_D)$ that is minimized by an isometry, we now consider algorithms for solving problems of the form
\begin{align}
\label{prog:ground_truth}
\arg \min_{ S \in \binom{[P]}{D}} \ell ( X_{. S})
\end{align}
where $\binom{[P]}{D} = \left\{ A \subseteq [P] : \left|A\right| = D \right\}$.
Brute force algorithms consider all possible solutions.
These algorithms are conceptually simple, but have the often prohibitive time complexity $O(C_\ell P^D)$ where $C_\ell$ is the cost of evaluating $\ell$.
Greedy algorithms consist of iteratively adding one element at a time to $ S$.
This algorithms have time complexity $O(C_\ell PD)$ and so are computationally more efficient than brute force algorithms, but can get stuck in local minima.
Formal definitions are given in Section \ref{sec:algorithms}.

Sometimes, it is possible to introduce an objective which convexifies problems of the above form  \citep{Liberti2004-iu, Abdi2013-th, Zhou2021-zz}.
Solutions
\begin{align}
\arg \min f(\beta) : Y  = X\beta 
\end{align}
to the overcomplete regression problem $Y = X \beta$ are a classic example \citep{Chen2001-hh}.
When $f(\beta) = \|\beta\|_0$, this problem is non-convex, and is thus suitable for greedy or brute algorithms, but when $f(\beta) =\|\beta\|_1$, the problem is convex, and may be solved efficiently via interior-point methods.
When the equality constraint is relaxed, Lagrangian duality may be used to reformulate as a so-called Lasso problem, which leads to an even richer set of optimization algorithms \citep{Hastie2016-cn}. % cite FISTA< glmnet, coordinate descent

The inspiration for this paper comes from an approach that stitches together pointwise conformal maps selected from a dictionary to form global embeddings \citep{Kohli2021-lr}.
That method constructs such isometries through greedy search, with putative dictionary features added one at a time.
